\documentclass[aspectratio=169,9pt]{beamer}
\usecolortheme{seahorse}
\setbeamertemplate{headline}{}
\setbeamertemplate{navigation symbols}{}
\setbeamertemplate{footline}{\leavevmode\hbox to \paperwidth{\hfill\raisebox{2pt}{\tiny\insertframenumber}\hspace{2mm}}\vskip1pt}
\usepackage{amsmath,amssymb,booktabs,tabularx,graphicx}
% Generated from ROOT-derived CSV sidecars on 2026-09-14.
\newcommand{\DeckDate}{14 September 2026}
\newcommand{\NCells}{21}
\newcommand{\NPerCell}{10,000}
\newcommand{\TotalEvents}{210,000}
\newcommand{\SeedPair}{26092601 / 8349041}
\newcommand{\COverN}{189.74\,mm/ns}
\newcommand{\AxialTime}{3.69\,ns}
\newcommand{\MeanPhotons}{19--20 thousand}

\newcommand{\VeffTwoHundred}{182.22\,mm/ns}
\newcommand{\VeffTwoOhFour}{181.65\,mm/ns}
\newcommand{\VeffTwoThirty}{181.61\,mm/ns}
\newcommand{\BetaTwoHundred}{16.19^{\circ}}
\newcommand{\BetaTwoOhFour}{16.79^{\circ}}
\newcommand{\BetaTwoThirty}{16.83^{\circ}}

\newcommand{\NpeTwoHundred}{528.4\pm1.3}
\newcommand{\NpeTwoOhFour}{398.1\pm1.0}
\newcommand{\NpeTwoThirty}{304.3\pm0.8}
\newcommand{\NpeTopTwoHundred}{5202.8\pm12.7}
\newcommand{\NpeTopTwoOhFour}{4767.8\pm11.8}
\newcommand{\NpeTopTwoThirty}{4086.2\pm10.0}
\newcommand{\NpeTotalTwoHundred}{6259.6\pm15.3}
\newcommand{\NpeTotalTwoOhFour}{5564.1\pm13.8}
\newcommand{\NpeTotalTwoThirty}{4694.7\pm11.5}
\newcommand{\EffTwoHundred}{32.542\%}
\newcommand{\EffTwoOhFour}{27.747\%}
\newcommand{\EffTwoThirty}{25.021\%}

\newcommand{\SigmaCenterTwoHundred}{76.49\pm0.86\,ps}
\newcommand{\SigmaCenterTwoOhFour}{71.47\pm0.80\,ps}
\newcommand{\SigmaCenterTwoThirty}{67.14\pm0.77\,ps}
\newcommand{\SigmaRangeTwoHundred}{76.4--82.0\,ps}
\newcommand{\SigmaRangeTwoOhFour}{71.5--84.3\,ps}
\newcommand{\SigmaRangeTwoThirty}{67.1--81.6\,ps}
\newcommand{\RmsCenterTwoHundred}{76.91\,ps}
\newcommand{\RmsCenterTwoOhFour}{72.44\,ps}
\newcommand{\RmsCenterTwoThirty}{67.77\,ps}

\newcommand{\SigmaXTwoHundred}{13.05\pm0.15\,mm}
\newcommand{\SigmaXTwoOhFour}{12.43\pm0.14\,mm}
\newcommand{\SigmaXTwoThirty}{11.42\pm0.13\,mm}
\newcommand{\SigmaXRangeTwoHundred}{13.0--14.6\,mm}
\newcommand{\SigmaXRangeTwoOhFour}{12.4--14.8\,mm}
\newcommand{\SigmaXRangeTwoThirty}{11.4--14.8\,mm}

\newcommand{\TZeroMeanTwoThirty}{4.0964\pm0.0008\,ns}
\newcommand{\TZeroChiTwoThirty}{160.4/137=1.17}
\newcommand{\TZeroRmsTwoThirty}{67.77\,ps}
\newcommand{\TZeroEdgeSigmaTwoThirty}{81.65\pm0.92\,ps}

\newcommand{\BestK}{3}
\newcommand{\BestKSigma}{66.00\pm0.76\,ps}
\newcommand{\FirstKSigma}{67.14\pm0.77\,ps}
\newcommand{\KImprovement}{1.14\pm1.08\,ps}
\newcommand{\BestMeanM}{5}
\newcommand{\BestMeanSigma}{57.59\pm0.66\,ps}
\newcommand{\MeanImprovement}{9.55\,ps}
\newcommand{\MeanImprovementPct}{14.2\%}

\newcommand{\EndOnlyReflectorGain}{1.09896\pm0.00852}
\newcommand{\EndTopReflectorGain}{1.01155\pm0.00830}
\newcommand{\EncounterRatio}{21.8}

\newcommand{\DecayTwoHundred}{2.1\,ns}
\newcommand{\DecayTwoOhFour}{1.8\,ns}
\newcommand{\DecayTwoThirty}{1.5\,ns}
\newcommand{\AttTwoHundred}{3.8\,m}
\newcommand{\AttTwoOhFour}{1.6\,m}
\newcommand{\AttTwoThirty}{1.2\,m}

\definecolor{endblue}{RGB}{30,90,180}
\definecolor{topgreen}{RGB}{20,145,80}
\definecolor{fastred}{RGB}{200,55,45}
\newcommand{\ps}{\,\mathrm{ps}}
\newcommand{\ns}{\,\mathrm{ns}}
\newcommand{\mm}{\,\mathrm{mm}}
\newcommand{\Source}[1]{\par\vfill{\scriptsize\color{gray}#1\par}}
\newcommand{\Result}[1]{{\Large\bfseries\color{endblue}#1}}

\title{What performance does the corrected optical model predict?}
\subtitle{Timing, light collection and longitudinal position in the current EndTop bar}
\author{René Ríos --- Universidad de La Serena}
\date{\DeckDate}

\begin{document}

\begin{frame}
\titlepage
\vspace{-6pt}
\centering\small Optical transport only: no SPTR, front-end, threshold or clock jitter.
\end{frame}

\begin{frame}{What are we trying to optimise?}
\begin{columns}[T]
\column{0.58\textwidth}
\includegraphics[width=\linewidth]{figures/geometry_current.pdf}
\column{0.38\textwidth}
\vspace{8pt}
\Result{Best optical time}\\[3pt]
with the fewest channels and a robust estimator.
\vspace{10pt}
\begin{tabular}{@{}ll@{}}
END & $8+8$ SiPMs\\
TOP & $70$ SiPMs\\
bar & $1400\times60\times10\,\mathrm{mm}^3$\\
scan & $x=0,\pm200,\pm500,\pm650\,\mathrm{mm}$
\end{tabular}
\vspace{9pt}

The comparison is between three plastics under one common geometry.
\end{columns}
\Source{Current corrected-transport grid; $N=\NPerCell$ per material and position.}
\end{frame}

\begin{frame}{From one muon to about 20,000 photons}
\begin{columns}[T]
\column{0.48\textwidth}
\centering
\vspace{9pt}
{\Large vertical 1 GeV $\mu^-$}\\[10pt]
\[
E_{\rm dep}\simeq 1.93\ \mathrm{MeV}
\]
\[
N_\gamma \simeq E_{\rm dep}\,Y
          \simeq (1.93)(9.7\mbox{--}10.4)\times10^3
\]
\Result{$N_\gamma\simeq 1.88$--$2.01\times10^4$}
\column{0.48\textwidth}
\vspace{8pt}
The detected population is shaped sequentially by
\begin{enumerate}
\item emission time and wavelength;
\item bulk absorption;
\item total internal reflection and wrapping;
\item active area and wavelength-dependent PDE.
\end{enumerate}
Only the earliest surviving photons determine the first-PE timing used here.
\end{columns}
\Source{Current data: mean $E_{\rm dep}=1.925$--$1.936$ MeV and $N_\gamma=18,781$--$20,070$ at $x=0$; event means.}
\end{frame}

\begin{frame}{How fast can information propagate?}
\begin{columns}[T]
\column{0.48\textwidth}
\textbf{Napkin prediction before Monte Carlo}
\[
v_{\max}=\frac{c}{n}=\COverN
\]
\[
t_{\min}(x=0)=\frac{700\mm}{189.74\mm/\ns}=\AxialTime
\]
This is the axial limit: $u_x=1$, no transverse motion.
\column{0.48\textwidth}
For a real ray,
\[
s=\frac{d}{|u_x|},\qquad P_{\rm bulk}=e^{-s/\lambda_{\rm att}}.
\]
Transverse components and reflections lengthen the path, so
\[
v_{\rm eff}<c/n.
\]
The current EJ-230 first-PE estimator gives $\langle T_0\rangle=\TZeroMeanTwoThirty$, about 0.41 ns above the axial limit.
\end{columns}
\Source{Analytical estimate uses $n=1.58$; measured time is current EJ-230, $x=0$, $N=\NPerCell$.}
\end{frame}

\begin{frame}{How much light reaches the END readout?}
\begin{columns}[T]
\column{0.64\textwidth}
\includegraphics[width=\linewidth]{figures/npe_vs_x.pdf}
\column{0.33\textwidth}
At $x=0$:
\begin{align*}
\text{EJ-200: }&\NpeTwoHundred\ \mathrm{pe/end}\\
\text{EJ-204: }&\NpeTwoOhFour\ \mathrm{pe/end}\\
\text{EJ-230: }&\NpeTwoThirty\ \mathrm{pe/end}
\end{align*}
Long attenuation wins the light budget: EJ-200 collects the most END light.

Near an end, the close arm gains strongly while the far arm remains photon-starved; the average yield rises but the timing balance worsens.
\end{columns}
\Source{Mean $(N_{pe,L}+N_{pe,R})/2\pm$ event-level SEM; current grid, $N=\NPerCell$/point.}
\end{frame}

\begin{frame}{The measured effective speed is about 182 mm/ns}
\begin{columns}[T]
\column{0.64\textwidth}
\includegraphics[width=\linewidth]{figures/delta_t_vs_x.pdf}
\column{0.33\textwidth}
\[
\Delta t\equiv t_R-t_L\simeq-\frac{2x}{v_{\rm eff}}
\]
\begin{tabular}{@{}lc@{}}
EJ-200 & $\VeffTwoHundred$\\
EJ-204 & $\VeffTwoOhFour$\\
EJ-230 & $\VeffTwoThirty$
\end{tabular}
\vspace{8pt}

This is 4.0--4.3\% below $c/n$. Interpreted as a single effective angle,
$\cos\beta_{\rm eff}=v_{\rm eff}/(c/n)$ gives $\beta_{\rm eff}=16.2$--$16.8^\circ$.
\end{columns}
\Source{ROOT linear fits to seven positions per material. The tiny statistical errors expose small non-linearity; $v_{\rm eff}$ is an effective slope, not a ray velocity.}
\end{frame}

\begin{frame}{Two times carry two different observables}
\begin{align*}
t_L &= t_0+\frac{L/2+x}{v_{\rm eff}}+\epsilon_L,
&t_R &= t_0+\frac{L/2-x}{v_{\rm eff}}+\epsilon_R.
\end{align*}
\begin{columns}[T]
\column{0.47\textwidth}
\centering
{\large\color{endblue}\textbf{Interaction time}}
\[
T_0\equiv\frac{t_L+t_R}{2}
\simeq t_0+\frac{L}{2v_{\rm eff}}.
\]
The first-order position dependence cancels.
\column{0.47\textwidth}
\centering
{\large\color{fastred}\textbf{Longitudinal position}}
\[
\Delta t\equiv t_R-t_L\simeq-\frac{2x}{v_{\rm eff}},
\quad
\sigma_x=\frac{v_{\rm eff}}{2}\sigma_{\Delta t}.
\]
The time difference is the position coordinate.
\end{columns}
\vspace{8pt}
Baseline here: $t_L$ and $t_R$ are the first detected photoelectron among the eight SiPMs on each END.
\Source{Definitions used event by event for every result that follows.}
\end{frame}

\begin{frame}{What does one timing distribution look like?}
\centering
\includegraphics[height=0.76\textheight]{figures/t0_fit_ej230_x0.pdf}
\vspace{-3pt}

The central Gaussian is physically adequate ($\chi^2/\mathrm{ndf}=\TZeroChiTwoThirty$); the global RMS, $\TZeroRmsTwoThirty$, is retained as a tail-sensitive cross-check.
\Source{EJ-230, $x=0$, $N=\NPerCell$; ROOT TF1 Gaussian refit over $\mu\pm2\sigma$; all entries shown in the left panel.}
\end{frame}

\begin{frame}{Optical timing is 67--84 ps across the bar}
\begin{columns}[T]
\column{0.65\textwidth}
\includegraphics[width=\linewidth]{figures/sigma_t0_vs_x.pdf}
\column{0.32\textwidth}
At the center:
\begin{align*}
\mathrm{EJ\!\!-200}:&\ \SigmaCenterTwoHundred\\
\mathrm{EJ\!\!-204}:&\ \SigmaCenterTwoOhFour\\
\mathrm{EJ\!\!-230}:&\ \SigmaCenterTwoThirty
\end{align*}
EJ-230 is fastest despite having the fewest PE.
\vspace{5pt}

Resolution worsens near the ends because $T_0$ averages one close arm with one long, weakly populated arm.
\vspace{5pt}

\textbf{Symmetry check:} every mirrored RMS and robust-width difference is below $2\sigma$; the visible residual wiggles are dominated by the Gaussian fit convention and finite statistics.
\end{columns}
\Source{Central Gaussian $\sigma(T_0)$ with ROOT fit uncertainty; current transport, first-PE END estimator, $N=\NPerCell$/point.}
\end{frame}

\begin{frame}{Why are the materials different?}
\begin{columns}[T]
\column{0.60\textwidth}
\includegraphics[width=\linewidth]{figures/material_comparison.pdf}
\column{0.37\textwidth}
Light yields are $10.0/10.4/9.7\times10^3$ photons/MeV; rise times are $0.9/0.7/0.5\ns$.
\vspace{4pt}

\textbf{EJ-200:} $\lambda_{att}=\AttTwoHundred$ preserves the most light, but $\tau_d=\DecayTwoHundred$ broadens early emission.
\vspace{5pt}

\textbf{EJ-230:} $\lambda_{att}=\AttTwoThirty$ yields the fewest PE, yet its $0.5\ns$ rise and $\tau_d=\DecayTwoThirty$ give the narrowest $T_0$ core.
\vspace{5pt}

The timing depends on the emph{early-photon} population, not the integrated PE count alone.
\end{columns}
\Source{Configured material constants; measured $N_{pe}$/end and central Gaussian timing at $x=0$ from the current grid.}
\end{frame}

\begin{frame}{Differential timing gives 11--15 mm longitudinal resolution}
\begin{columns}[T]
\column{0.65\textwidth}
\includegraphics[width=\linewidth]{figures/sigma_x_vs_x.pdf}
\column{0.32\textwidth}
\[
\sigma_x=\frac{v_{\rm eff}}{2}\sigma(\Delta t)
\]
At $x=0$:
\begin{align*}
\mathrm{EJ\!\!-200}:&\ \SigmaXTwoHundred\\
\mathrm{EJ\!\!-204}:&\ \SigmaXTwoOhFour\\
\mathrm{EJ\!\!-230}:&\ \SigmaXTwoThirty
\end{align*}
The best central point is EJ-230; all materials approach 15 mm near an end.
\end{columns}
\Source{Central ROOT Gaussian width of $\Delta t$; uncertainty propagates the $\sigma_{\Delta t}$ and slope-fit errors.}
\end{frame}

\begin{frame}{More light is not necessarily better timing}
\begin{columns}[T]
\column{0.58\textwidth}
\includegraphics[width=\linewidth]{figures/npe_vs_sigma.pdf}
\column{0.39\textwidth}
Two conceptual populations contribute:
\begin{itemize}
\item \textbf{TIR-guided photons}: early, shorter optical paths;
\item \textbf{reflector-recovered photons}: rescued light, often after more encounters and longer paths.
\end{itemize}
Current post-repair reflector A/B:
\begin{align*}
G_{\rm END-only}&=\EndOnlyReflectorGain,\\
G_{\rm EndTop}&=\EndTopReflectorGain.
\end{align*}
The reflector adds far more END light when the bar is fully wrapped; it does not guarantee a narrower early-time distribution.
\end{columns}
\Source{Plot: all 21 current cells. Reflector A/B: EJ-204, $x=0$, $N=2000$/arm, corrected transport; approximate independent-arm SEM.}
\end{frame}

\begin{frame}{Why should the first photon be optimal?}
\begin{columns}[T]
\column{0.64\textwidth}
\includegraphics[width=\linewidth]{figures/order_scan.pdf}
\column{0.33\textwidth}
\[
T_0(k)=\frac{t_L^{(k)}+t_R^{(k)}}{2}
\]
The optimum is
\[
m=\BestMeanM:\quad \BestMeanSigma.
\]
Against the first PE ($\FirstKSigma$), the gain is $\MeanImprovement$
($\MeanImprovementPct$). The k-th photon alone has only a shallow minimum at
$k=\BestK$: $\BestKSigma$.
\vspace{4pt}

Averaging the first few arrivals suppresses first-photon fluctuations; including too many eventually admits the broader late population.
\end{columns}
\Source{EJ-230, $x=0$, $N=\NPerCell$; k-th and arithmetic mean-of-first-m scans on the same events; no position-dependent selection.}
\end{frame}

\begin{frame}{TOP contains early timing information---but at a high channel cost}
\begin{columns}[T]
\column{0.59\textwidth}
\includegraphics[width=\linewidth]{figures/end_vs_top.pdf}
\column{0.38\textwidth}
At $x=0$, TOP detects $4.1$--$5.2\times10^3$ PE across 70 sensors, versus 304--528 PE per eight-sensor END.
\vspace{5pt}

The plotted 2.9--3.1 ps TOP width is the 	extbf{global first photon among all 70 channels}: an optical lower bound with a position-local sensor.
\vspace{5pt}

No validated threshold, pulse shape or channel-combination estimator exists yet. These data do not establish an END+TOP detector resolution.
\end{columns}
\Source{Current $x=0$ cells, $N=\NPerCell$; END and TOP observables are deliberately kept separate.}
\end{frame}

\begin{frame}{Current optical performance}
\centering
\small Common configuration: EndTop, $16+70$ SiPMs, $x=0$ unless stated, first-PE END estimator.
\vspace{7pt}
\begin{tabular}{@{}lccc@{}}
\toprule
 & EJ-200 & EJ-204 & \textbf{EJ-230}\\
\midrule
$N_{pe}$/end at center & $\NpeTwoHundred$ & $\NpeTwoOhFour$ & $\NpeTwoThirty$\\
$v_{\rm eff}$ & $\VeffTwoHundred$ & $\VeffTwoOhFour$ & $\VeffTwoThirty$\\
$\sigma_{T_0}$ at center & $\SigmaCenterTwoHundred$ & $\SigmaCenterTwoOhFour$ & \textbf{$\SigmaCenterTwoThirty$}\\
$\sigma_{T_0}$ over x scan & $\SigmaRangeTwoHundred$ & $\SigmaRangeTwoOhFour$ & $\SigmaRangeTwoThirty$\\
$\sigma_x$ at center & $\SigmaXTwoHundred$ & $\SigmaXTwoOhFour$ & \textbf{$\SigmaXTwoThirty$}\\
\bottomrule
\end{tabular}
\vspace{10pt}

\Result{EJ-230 gives the best intrinsic END timing and position resolution.}

EJ-200 remains the light-yield winner. These are optical widths, not experimental detector resolutions.
\Source{All entries recomputed from current corrected-transport ROOTs; $N=\NPerCell$/cell; uncertainties as defined on preceding slides.}
\end{frame}

\begin{frame}{What should we optimise next?}
\begin{columns}[T]
\column{0.48\textwidth}
\textbf{Physics already measured}
\begin{itemize}
\item END first-PE baseline: 67--84 ps;
\item the mean of the first five END PE improves 67.1 to 57.6 ps;
\item TOP carries very early local information;
\item attenuation controls light, decay time controls early timing.
\end{itemize}
\column{0.48\textwidth}
\textbf{Design question still open}
\begin{itemize}
\item build one electronics-aware TOP estimator;
\item compare fixed channel counts and thresholds;
\item validate SPTR, pulse shape and TDC together;
\item optimise $\sigma_t$ against channel count and complexity.
\end{itemize}
\end{columns}
\vspace{9pt}
\centering
\Result{The next decision is not ``more sensors?'' but ``how much validated timing per channel?''}
\Source{No electronics-convolved or experimental resolution is inferred in this presentation.}
\end{frame}

\appendix
\begin{frame}{Backup: dataset and common configuration}
\small
\begin{tabularx}{\linewidth}{@{}lX@{}}
\toprule
Cells & 3 materials $\times$ 7 positions = \NCells; \NPerCell{} generated events per cell\\
Geometry & $1400\times60\times10\mm^3$ EndTop bar; 8 SiPMs on each END; 70 TOP SiPMs\\
Beam & vertical 1 GeV $\mu^-$; $x=0,\pm200,\pm500,\pm650\mm$\\
Timing content & individual detected-PE time, SiPM and face; event-level yields and energy deposition\\
Electronics & SPTR = 0; no pulse formation, threshold, TDC or clock jitter\\
Statistics & 210,000 events; all cells readable and complete\\
\bottomrule
\end{tabularx}
\vspace{7pt}
Each input was checked against its recorded geometry, material, seeds, event count and binary before analysis. No new simulation was required.
\Source{Technical dataset audit is retained in the v9 source directory.}
\end{frame}

\begin{frame}{Backup: exact observable and fit convention}
\small
For each event, pool the detected photoelectrons over the eight SiPMs at each END and sort by arrival time:
\[
t_L^{(k)},\quad t_R^{(k)},\qquad
T_0(k)=\tfrac12[t_L^{(k)}+t_R^{(k)}],\quad
\Delta t(k)=t_R^{(k)}-t_L^{(k)}.
\]
Baseline: $k=1$. Resolution: central Gaussian $\sigma$ from a ROOT TF1 fit. The first pass uses the median $\pm2(q_{84}-q_{16})/2$; the reported refit uses $\mu\pm2\sigma$.
\vspace{5pt}

The global RMS is calculated on all events. The robust cross-check is
$q_w=(q_{84}-q_{16})/2$. RMS and $q_w$ errors use 300 deterministic event-bootstrap replicas; Gaussian errors come from the ROOT fit.
\vspace{5pt}

The construction is explicitly symmetric: global IDs 0--7 form left END and 8--15 right END, then $T_0$ averages the two arms with equal weight.
\end{frame}

\begin{frame}{Backup: all seven EJ-230 fits}
\centering
\includegraphics[height=0.79\textheight]{figures/fit_grid_ej230.pdf}
\Source{Exact histograms and stored ROOT TF1 objects used by the quoted resolution; $N=\NPerCell$ each.}
\end{frame}

\begin{frame}{Backup: EJ-230 fit diagnostics}
\scriptsize
\begin{tabular}{@{}rcccccc@{}}
\toprule
$x$ [mm] & $\sigma_G$ [ps] & fit mean [ns] & $\chi^2$/ndf & fit window [ns] & RMS [ps] & $q_w$ [ps]\\
\midrule
-650 & $81.27\pm0.92$ & 4.0912 & 1.65 & [3.929,4.255] & 85.05 & 84.98\\
-500 & $76.13\pm0.87$ & 4.0981 & 1.24 & [3.946,4.252] & 79.22 & 78.51\\
-200 & $69.38\pm0.81$ & 4.0968 & 1.15 & [3.959,4.235] & 69.21 & 69.33\\
0 & $67.14\pm0.77$ & 4.0964 & 1.17 & [3.962,4.231] & 67.77 & 66.95\\
+200 & $69.32\pm0.78$ & 4.0989 & 1.34 & [3.960,4.238] & 70.11 & 70.22\\
+500 & $78.04\pm0.89$ & 4.1014 & 1.58 & [3.945,4.259] & 80.51 & 80.67\\
+650 & $81.65\pm0.92$ & 4.0929 & 1.96 & [3.930,4.257] & 85.50 & 86.29\\
\bottomrule
\end{tabular}
\vspace{8pt}

All 21 material/position records, including fit status and both width uncertainties, are stored in the numerical sidecar. Every ROOT fit status is zero.
\Source{$q_w=(q_{84}-q_{16})/2$; RMS and $q_w$ errors are omitted from this compact table and retained in the sidecar.}
\end{frame}

\begin{frame}{Backup: mirrored EJ-230 shapes}
\centering
\includegraphics[width=0.98\linewidth]{figures/mirror_overlays_ej230.pdf}

After centering each distribution at its own mean, the mirrored shapes overlap. All panels use the same $[-0.30,+0.30]$ ns range, 120 bins and unit-area normalization.
\Source{Current EJ-230 ROOT events; $N=\NPerCell$ per position. No refit or width rescaling is applied.}
\end{frame}

\begin{frame}{Backup: the width definition controls the visible wiggles}
\begin{columns}[T]
\column{0.64\textwidth}
\includegraphics[width=\linewidth]{figures/width_estimators_ej230.pdf}
\column{0.33\textwidth}
For EJ-230, every mirrored difference is below $1.9\sigma$ for all three definitions.
\vspace{6pt}

Across all nine mirrored pairs:
\begin{itemize}
\item RMS: $|z|<1.94$;
\item robust $q_w$: $|z|<1.90$;
\item Gaussian: two EJ-200 pairs reach $2.14\sigma$ and $3.06\sigma$.
\end{itemize}
The robust definitions support left--right symmetry; the largest visible Gaussian deviations are fit sensitivity plus finite statistics.
\end{columns}
\Source{RMS/$q_w$: 300 event-bootstrap replicas; Gaussian: ROOT fit covariance. Differences are $+x$ minus $-x$.}
\end{frame}

\begin{frame}{Backup: mirrored-width differences for all materials}
\fontsize{5.8}{7}\selectfont
\begin{tabular}{@{}lrrrrrrr@{}}
\toprule
 & $|x|$ & $\Delta\sigma_G$ [ps] & $z_G$ & $\Delta$RMS [ps] & $z_R$ & $\Delta q_w$ [ps] & $z_q$\\
\midrule
EJ-200 & 200 & $+3.89\pm1.27$ & +3.06 & $+1.51\pm0.78$ & +1.94 & $+1.77\pm1.05$ & +1.69\\
EJ-200 & 500 & $-0.64\pm1.30$ & -0.49 & $+0.48\pm0.77$ & +0.63 & $+0.82\pm0.99$ & +0.83\\
EJ-200 & 650 & $+2.72\pm1.27$ & +2.14 & $+0.61\pm0.79$ & +0.76 & $+0.33\pm1.06$ & +0.31\\
EJ-204 & 200 & $+0.08\pm1.16$ & +0.07 & $-0.44\pm0.71$ & -0.61 & $+0.29\pm0.94$ & +0.31\\
EJ-204 & 500 & $+1.30\pm1.29$ & +1.00 & $-0.09\pm0.82$ & -0.11 & $+1.27\pm1.10$ & +1.15\\
EJ-204 & 650 & $-1.18\pm1.34$ & -0.88 & $+0.25\pm0.89$ & +0.28 & $-0.44\pm1.20$ & -0.37\\
EJ-230 & 200 & $-0.06\pm1.12$ & -0.05 & $+0.89\pm0.71$ & +1.26 & $+0.90\pm0.95$ & +0.94\\
EJ-230 & 500 & $+1.91\pm1.24$ & +1.54 & $+1.29\pm0.87$ & +1.48 & $+2.16\pm1.14$ & +1.89\\
EJ-230 & 650 & $+0.38\pm1.30$ & +0.29 & $+0.45\pm0.92$ & +0.49 & $+1.31\pm1.26$ & +1.04\\
\bottomrule
\end{tabular}
\vspace{5pt}

No RMS or robust-width difference reaches $2\sigma$. There is no coherent sign across material or position.
\Source{$\Delta w=w(+|x|)-w(-|x|)$. Bootstrap seed series begins at 26091443.}
\end{frame}

\begin{frame}{Backup: scope after the symmetry check}
\small
\begin{tabularx}{\linewidth}{@{}lX@{}}
\toprule
Reproduced & Optical END $T_0$ and $\Delta t$; k-th and mean-first-m scans; Gaussian, RMS and robust widths; mirrored-pair diagnostics.\\
\midrule
Construction check & left IDs 0--7, right IDs 8--15; equal-weight $T_0$; zero hit-count mismatches in 210,000 events; mirrored yields and mean times swap within finite-sample fluctuations.\\
\midrule
Unavailable & validated END+TOP estimator; electronics-convolved or experimental resolution; full per-track TIR/recovered timing split; physically validated weighted channel mean.\\
\bottomrule
\end{tabularx}
\vspace{6pt}
Historical v6/v7 timing values remain excluded. The current main curve keeps the central Gaussian convention, with the robust symmetry cross-check stated explicitly.
\Source{Exact software provenance remains in the v8 validation deck and the v9 technical report; no new simulation was run.}
\end{frame}


\end{document}
